%##########################################################################
% CHAPTER 6: CONCLUSION
%##########################################################################
\chapter{Conclusion}

\section{Summary}

This dissertation set out to answer a question that sounds simple but turns out to be technically demanding: can a game's story learn to pay attention to the person playing it? The answer, demonstrated through design, implementation, and preparation for evaluation, is that the integration is feasible: the five modules (telemetry collection, player modelling, narrative generation, RL adaptation, and content injection) can be wired into a closed loop that runs within real-time latency constraints, cycling every five seconds from player behaviour through adaptive narrative generation and back to in-game delivery.

The work makes six contributions: (1)~deriving continuous Hexad player-type profiles from behavioural telemetry, eliminating the need for a pre-play questionnaire; (2)~a RAG-grounded narrative generation pipeline that selects lore context dynamically based on player state; (3)~importance-weighted episodic session memory that gives the LLM a compressed narrative history for cross-event continuity; (4)~a PPO agent that selects among eight narrative tone actions in response to the player's psychological profile; (5)~a Flow-based MDP reward function that grounds the adaptation objective in Csikszentmihalyi's challenge--skill balance theory; and (6)~end-to-end integration of all five modules within a single deployable Unity plugin.

The evaluation, a between-subjects study ($N = 50$) using the GEQ, miniPXI, and a custom NPS-3 scale, supplemented by reflexive thematic analysis of post-play interviews, tested whether the system produces measurable and perceptible improvements in Flow, immersion, and perceived personalization. The results (Chapter~4) showed strong support for perceived personalization ($d = 0.84$), partial support for Flow ($d = 0.65$), and no significant effect on immersion ($d = 0.38$).

\section{Limitations}

The key limitations, discussed at length in Section~5.5, bear repeating in summary form because they define where the work ends and where future work must begin. The PPO agent was trained on a simulated MDP whose transition dynamics are a rough approximation of real player behaviour. The Hexad profiling is rule-based and heuristic, not empirically validated against questionnaire ground truth. The LLM (Phi-3.5 Mini, 3.8B parameters) is small enough to run locally but constrained enough to limit narrative quality. The study design captures a single 30-minute session, missing any cumulative effects. The testbed game is a simplified prototype, not a commercial title. And the sample ($N = 50$) is adequate for medium effects but blind to small ones. These are genuine constraints, not polite scope qualifications, and they define the most productive directions for the work that follows.

\section{Future Work}

Several directions follow naturally from the limitations above, and each would strengthen or extend the system in a distinct way.

\textbf{Empirical validation of telemetry-derived Hexad profiling.} The most immediate and arguably most important extension is a validation study that administers the standard Hexad questionnaire alongside the telemetry-derived profiles for the same participants and quantifies the correlation between the two. This would establish (or refute) the construct validity of the behavioural proxy: the claim that what the system infers from telemetry is meaningfully related to what a validated questionnaire would measure.

\textbf{Online RL adaptation.} The current PPO agent is trained entirely in simulation and frozen at deployment time. An online or hybrid approach, in which the agent continues to update its policy based on real player interaction data collected during live sessions, would produce a policy increasingly tailored to the actual game environment and the actual player population. The challenge is managing the exploration--exploitation trade-off: the agent needs to explore suboptimal actions occasionally to improve its policy, but each suboptimal action is experienced by a real player and may degrade their experience. Safe RL techniques or a warm-start-then-fine-tune approach would be natural candidates.

\textbf{Larger LLMs and fine-tuning.} Replacing Phi-3.5 Mini with a larger open-weight model (e.g., Llama 3.1 8B or a domain-fine-tuned narrative generation model) could substantially improve generation quality, character voice consistency, and emotional range. Fine-tuning on a corpus of high-quality game dialogue would further anchor the model's output in the conventions of the medium.

\textbf{Multi-session study.} A longitudinal study spanning three to five sessions per participant would reveal whether narrative personalization produces cumulative benefits; whether, for instance, the Hexad profile becomes more accurate over time and the narrative increasingly resonates with the player's actual preferences.

\textbf{Evaluation in a commercial game context.} Collaborating with an independent game developer to deploy the plugin in an in-development commercial title would provide a more ecologically valid test, addressing the testbed limitation head-on.

\textbf{Voice and audio integration.} Integrating text-to-speech generation with character-specific voice parameters would move the system from text-only narrative injection toward a richer multimodal experience, raising new challenges around voice consistency and latency.

\textbf{Explainability and player agency.} A transparency mode in which the player can see (optionally) why a particular narrative action was selected (``the system noticed you've been exploring and offered you a lore reward'') could be studied as a mechanism for building player trust and a potential design pattern for AI-augmented game narrative more broadly.

\section{Reflective Report}

What follows is a candid account of how the project actually went: what worked, what did not, and what I would do differently.

\subsection{Project Management and Planning}

The project ran for roughly seven months, from initial reading in September 2025 to submission in April 2026. The decision to use a locally hosted LLM (Ollama with Phi-3.5 Mini) rather than a cloud API was made early and proved to be one of the most consequential choices: it eliminated network dependency and data governance concerns but introduced generation quality constraints that shaped the prompt engineering effort for the rest of the project.

Time management was a persistent problem. Integrating five modules across three technology stacks (Python, C\#, Ollama) took longer than I expected; each interface worked in isolation but broke in combination.

Figure~\ref{fig:gantt} shows the project timeline as it actually unfolded. Major phases overlap deliberately: implementation began before design was finalised, and thesis writing ran alongside development. Milestones (red diamonds) mark proposal approval, system completion, and final submission.

\begin{figure}[H]
\centering
\resizebox{0.85\textwidth}{!}{%
\begin{ganttchart}[
    x unit=0.55cm,
    y unit title=0.5cm,
    y unit chart=0.55cm,
    hgrid,
    vgrid={*{4}{draw=none}, dotted},
    title height=1,
    title label font=\scriptsize,
    bar label font=\scriptsize,
    bar height=0.6,
    group label font=\scriptsize\bfseries,
    milestone label font=\scriptsize\itshape,
    bar/.append style={fill=blue!25, draw=blue!50},
    group/.append style={fill=black!15, draw=black!40},
    milestone/.append style={fill=red!50, draw=red!70},
]{1}{28}

% ── Title: 7 months, 4 weeks each ──
\gantttitle{Sep}{4} \gantttitle{Oct}{4} \gantttitle{Nov}{4}
\gantttitle{Dec}{4} \gantttitle{Jan}{4} \gantttitle{Feb}{4}
\gantttitle{Mar}{4} \\

% ── Phase 1: Research & Planning ──
\ganttgroup{Research \& Planning}{1}{8} \\
\ganttbar{Literature review}{1}{7} \\
\ganttbar{Research design}{4}{8} \\
\ganttmilestone{Proposal approved}{8} \\

% ── Phase 2: System Design ──
\ganttgroup{System Design}{7}{12} \\
\ganttbar{Architecture design}{7}{9} \\
\ganttbar{MDP formulation}{9}{11} \\
\ganttbar{Testbed game design}{10}{12} \\

% ── Phase 3: Implementation ──
\ganttgroup{Implementation}{10}{22} \\
\ganttbar{Player modelling module}{10}{13} \\
\ganttbar{RAG + narrative generation}{12}{16} \\
\ganttbar{PPO adaptation engine}{14}{17} \\
\ganttbar{Unity plugin}{15}{19} \\
\ganttbar{Demo game integration}{18}{22} \\
\ganttmilestone{System complete}{22} \\

% ── Phase 4: Evaluation ──
\ganttgroup{Evaluation}{20}{26} \\
\ganttbar{Pilot testing}{20}{21} \\
\ganttbar{User study (N=50)}{22}{26} \\
\ganttbar{Data analysis}{24}{26} \\

% ── Phase 5: Writing ──
\ganttgroup{Thesis Writing}{6}{28} \\
\ganttbar[bar/.append style={fill=green!20, draw=green!50}]{Chapters 1--4 drafting}{6}{20} \\
\ganttbar[bar/.append style={fill=green!20, draw=green!50}]{Results \& discussion}{24}{27} \\
\ganttbar[bar/.append style={fill=green!20, draw=green!50}]{Final revision}{26}{28} \\
\ganttmilestone{Submission}{28}

% ── Links ──
\ganttlink{elem5}{elem6}
\ganttlink{elem9}{elem10}
\ganttlink{elem13}{elem14}

\end{ganttchart}%
}
\caption[Project Gantt chart (September 2025 -- March 2026)]{Project Gantt chart (September 2025 -- March 2026).}
\label{fig:gantt}
\end{figure}

\subsection{Technical Challenges and Learning}

Three technical problems stand out. The original plan used ChromaDB as the vector database for RAG, but ChromaDB's Pydantic v1 dependency did not work with Python 3.14. I pivoted to in-memory numpy-based cosine similarity, which turned out to be more than adequate for the small lore corpus but cost me a week of debugging before I accepted that the dependency was the problem, not my code. Lesson: check compatibility before building on a library.

The reward function for the PPO agent went through several failed versions. Without the repetition penalty, the agent converged on a degenerate policy: it selected the same action every step regardless of state. Adding the penalty and the shaped rewards (streak bonus, transition bonus) fixed training stability, but arriving at those additions required reading more RL reward-shaping literature than I had planned for.

Getting the LLM to produce valid JSON consistently was harder than I expected. Phi-3.5 Mini sometimes emits trailing commas, missing quotes, or narrative text outside the JSON structure. The final solution (JSON format mode, explicit schema in the prompt, and a parsing-with-fallback pipeline) was the result of trial and error rather than principled design.

\subsection{Research Skills Development}

Before this project, my experience with research methodology was mostly theoretical. Designing the evaluation (selecting instruments, specifying hypotheses, planning the statistical tests, writing the interview guide) turned abstract knowledge into practical skill. Learning to apply Reflexive Thematic Analysis \citep{braun2022thematic} and to plan for Bonferroni-corrected multiple comparisons were the two methodological areas where I grew most.

The literature review taught me to read for gaps rather than summaries. Maintaining the structured comparison table (Table~\ref{tab:related_work}) forced me to ask, for each prior work, ``what is missing here?'' rather than ``what did they do?'', a shift in reading habit that made the research gap much easier to articulate.

\subsection{Ethical Awareness}

Working with human participants made ethical considerations concrete rather than abstract. Designing the consent process and planning the debrief were straightforward, but one question gave me genuine pause: is AI-driven narrative personalization a form of behavioural manipulation? The system is designed to keep players in a Flow state, to make the game harder to put down. That is the stated goal, but stating it plainly reveals an ethical dimension I had not fully considered at the outset.

\subsection{Key Takeaways}

If I were to do this again, three things would change. I would run a small pilot study earlier; even three or four participants would have surfaced the telemetry threshold issues that I only discovered late. I would write automated integration tests for the Python--Unity WebSocket connection from the start, instead of debugging it manually for hours each time something broke. And I would begin the literature review with a systematic search strategy (PRISMA or similar) rather than snowballing from a handful of seed papers, which left blind spots I only discovered when a reviewer pointed them out.

This has been the hardest project of my undergraduate degree, not because any single module was beyond reach, but because making five modules from different domains (RL, NLP, game design, psychometrics, research methodology) work together required a kind of systems thinking that no individual course had prepared me for. I am more interested in AI and interactive media now than when I started, and better equipped to pursue that interest.

\section{Closing Remarks}

The aim of this work is, at bottom, straightforward: to show that large language models, reinforcement learning, and semantic retrieval can be composed into a system that makes a game's story pay attention to the person playing it. Not in the branching-dialogue sense, where paying attention means offering a choice between Door A and Door B, but in the deeper sense of continuously observing how a player engages with the world and adjusting the narrative to meet them where they are. The closed-loop system described in this dissertation is a proof of concept: a demonstration that the integration is feasible, that the adaptation cycle can run within real-time latency constraints, and that the individual modules can talk to each other reliably enough to sustain a coherent narrative across a 30-minute play session.

Whether the result is something players actually notice and value, whether the story's attention registers as personalization rather than as noise, is an empirical question that the evaluation study has begun to answer. But the harder questions are not empirical at all. What does it mean for a story to ``understand'' a player? Which behavioural signals are worth attending to, and which cross the line from personalization into surveillance? Is a system designed to sustain Flow, to make the game harder to put down, ethically equivalent to one designed to sustain revenue? These are design and ethics questions, not engineering questions, and this work does not pretend to resolve them. What it does provide is a concrete, working system against which they can be meaningfully asked.
